% =====================================================================
%  SUBSECTION: The Yang--Mills mass gap
%  Drafted 2026-07-03 per JS: having committed to all mass as
%  associator-geodesic length (including the proton's), the paper must
%  say why it does not extend the claim to the Clay Millennium problem
%  -- "it would seem strange to get so far and not make the extra
%  effort" -- with an explanation and a forward-looking suggestion.
%
%  PLACEMENT: into confinement_section_2026-06-13.tex, immediately
%  AFTER the "All mass is geodesic length" subsection (i.e. after its
%  closing gravity/E=mc^2 forward sentence, the paragraph beginning
%  "One forward-looking consequence deserves a single sentence...")
%  and BEFORE \subsection{Confinement and pre-spatial co-location}.
%
%  REGISTER (deliberate): this is the highest-crankery-risk topic in
%  physics; the subsection names the problem precisely so that no claim
%  on it can be read into the section by implication, declines it with
%  specific reasons (same kinematic/dynamical boundary as the QED
%  scope statement), then records the framework-native structural seed
%  (conditional, flagged) and two forward directions. NO new open
%  problem is minted: the framework's pieces of the gap are already
%  op:associator-geodesic, op:glueball-spectrum, op:qcd-scale,
%  op:confinement-potential, and the four translation interfaces are
%  given as a prose programme.
%
%  DEPENDENCY: references op:qcd-scale, which is DEFINED in
%  associator_debt_higgs_mechanism_2026-07-03.tex -- paste that
%  subsection into the boson-masses section first, or this reference
%  dangles. (Fallback if you prefer: drop the op:qcd-scale clause.)
%
%  CROSS-REFS (all existing): sec:interactions, sec:space-of-computations,
%  sec:masses, op:associator-geodesic, op:glueball-spectrum,
%  op:confinement-potential, op:qcd-scale [see dependency].
%  NEW citation keys (RIS in mass_gap_refs_2026-07-03.ris):
%  jaffewitten2006, cubitt2015spectralgap.
%
%  OPTIONAL one-line addition to "What is settled and what is owed",
%  end of the owed paragraph:
%    "The Yang--Mills mass gap is named in \S\ref{sec:mass-gap}
%     precisely so that no claim on it is read into this section by
%     implication."
% =====================================================================

\subsection{The Yang--Mills mass gap}
\label{sec:mass-gap}

A reader who has followed this section to here will have noticed that its commitments walk up to
the edge of a famous problem, and it would be evasive not to name it. The Clay Millennium
formulation asks for two things \autocite{jaffewitten2006}: a mathematically rigorous construction
of quantum Yang--Mills theory on $\mathbb{R}^{4}$, for any compact simple gauge group, meeting the
Wightman axioms or standards of comparable stringency; and a proof that the constructed theory has
a mass gap $\Delta > 0$ --- that its lightest excitation is strictly massive. As physics the gap is
not in doubt: lattice computation exhibits it, and its lightest carrier is the scalar glueball near
$1.7$ GeV whose derivation this section has already recorded as
Open Problem~\ref{op:glueball-spectrum}. The Millennium problem is a problem of proof, not of
phenomenology. This section has asserted that confinement is exact and structural, that glueballs
are closed colour networks, and that all mass --- theirs included --- is associator-geodesic length.
Having gone that far, we owe the reader a plain statement of why we go no further, and of what
going further would require.

The reasons are the boundary already drawn for quantum electrodynamics in \S\ref{sec:interactions},
here made specific. The mass gap is a statement about the spectrum of a quantum dynamics, and this
framework has not constructed that dynamics: there is no Hamiltonian here whose spectrum could be
bounded, no Hilbert space of the field theory, and no continuum $\mathbb{R}^{4}$ to put one on ---
spacetime itself being emergent in this framework, its construction deferred
(\S\ref{sec:space-of-computations}). One cannot prove a spectral gap of an operator one has not
built. Beneath that lies the framework's own open centre: the map from network geometry to physical
mass is Open Problem~\ref{op:associator-geodesic}, and the statement ``every closed colour network
has strictly positive geodesic length'' cannot be a theorem before the functional it concerns
exists. Even were both supplied, the translation into the Millennium problem's terms crosses four
interfaces, each a research programme in its own right: the continuum construction of the emergent
spacetime; the translation of the categorical and entanglement-theoretic statements of this paper
into an axiomatic form; the passage from the algebraic obstruction of this section to the standard
criteria of confinement, the Wilson-loop area law above all; and the spectral proof itself. We claim
none of these. One further limit of scope deserves its own sentence: the Millennium problem is posed
for every compact simple gauge group, while the structural story of this section is specific to the
colour sector as it arises here, on the octonionic register --- our route, even if completed, would
speak to $SU(3)$, not to the problem in its full generality.

What the framework does supply, and what we record because it gives the gap a concrete identity, is
a structural reason to expect one. The photon is massless because the electroweak sector contains a
debt-free direction: the associator-blind $\mathbb{C}$, along which nothing is owed
(\S\ref{sec:masses}). The colour sector contains no such direction --- colour lives on the register
exactly where the associator is non-trivial --- so no closed colour network can be debt-free. And
closed networks are topologically discrete: the minimal closure is a finite structure, and there is
no continuum of ever-smaller closed networks by which geodesic lengths could accumulate towards
zero. If, as the framework's mass formulas everywhere suggest, the geodesic length of a closed
network is fixed by its discrete topology rather than by any continuously variable modulus, then the
spectrum of pure-colour states is bounded below by the length of the minimal closed network. On this
reading the gap exists because the colour sector has no massless direction and no arbitrarily small
closed structure; and $\Delta$ is then not an abstract bound but a particular number --- the mass of
the lightest glueball, precisely the object of Open Problem~\ref{op:glueball-spectrum}, with its
scale in physical units the object of Open Problem~\ref{op:qcd-scale}. Everything in this paragraph
is conditional on the open map: it is a heuristic the framework makes natural, not a theorem the
framework possesses, and we flag it as such.

Two forward directions are worth leaving on the table. The first reframes the problem's notorious
difficulty. Cubitt, P\'{e}rez-Garc\'{i}a and Wolf proved that the spectral-gap question is, in
general, undecidable: for translationally invariant lattice Hamiltonians it encodes the halting
problem, and their own statement of the result names the Yang--Mills gap conjecture as a particular
case of the general problem so shown undecidable \autocite{cubitt2015spectralgap}. For a framework
whose spine is self-referential non-computability this is orientation rather than discouragement.
It says that any proof must exploit special structure --- no general algorithmic route exists ---
and it locates the difficulty exactly where the bump-in-the-carpet principle of this framework
expects to find it. The suggestion we record, in the same register as the other prospects of this
paper --- a direction with a possible path, not a claim --- is that the special structure is the
division-algebra rigidity itself, the Hurwitz--Adams rigidity that terminated the Hopf tower, and
that a proof along this route would show the geodesic functional on closed colour networks to be
bounded below by the minimal-network length, the discreteness of network topology doing the work
that lattice regularisation does elsewhere. The second direction is nearer, and does not wait on the
emergent-distance machinery: in a fixed lattice geometry, the question of whether the
budget-saturation constraint of this section reproduces area-law Wilson loops is one that
tensor-network methods can address now, and it is the interface at which this framework's account of
confinement first becomes commensurable with the standard one. We leave the problem, and the route,
in the record of open problems rather than in the claims of this paper.
